% HD-I22. GPT-6 Astra Pro /root/imprimitive, 21 November 2026.
% Inherited effort; numerical setting not exposed.
\section{HD-I22: the exact field of the remaining rank-three coefficient group}
\label{sec:rank-three-exact-field}

Let
\[
 \mathcal K=\{\mathbf Q(i),\mathbf Q(\sqrt{-2}),
 \mathbf Q(\sqrt{-7}),\mathbf Q(\sqrt{-11}),
 \mathbf Q(\sqrt{-19}),\mathbf Q(\sqrt{-43}),
 \mathbf Q(\sqrt{-67}),\mathbf Q(\sqrt{-163})\}.
                                                               \tag{I22.1}
\]
For a finite absolutely irreducible subgroup
$\Gamma\subset\operatorname{GL}_3(K)$, $K\in\mathcal K$,
whose exact character field is $K$, the following list is
complete. Over $\mathbf Q(i)$, $\Gamma$ has an actual
$K$-defined system of three lines. Over
$K_7=\mathbf Q(\sqrt{-7})$, it is $K$-linearly conjugate
to one of
\[
 C_7\rtimes C_3,\quad
 \{\pm I\}\times(C_7\rtimes C_3),\quad
 L,\quad \{\pm I\}\times L,\qquad L\simeq\operatorname{PSL}_2(7),
                                                               \tag{I22.2}
\]
in the actual degree-three models below. For the other six
fields in (I22.1) there is no such group of exact character
field $K$.

Applied through the actual curve and corner maps of HD-S45,
the first two groups in (I22.2) cannot give period rank nine.
Consequently the only remaining multiplicity-three tensor
source from that reduction is
\[
 K=K_7,\qquad H=S_s\times\Gamma,\qquad
 \Gamma=L\ \hbox{or}\ \{\pm I\}\times L,
 \qquad s\ge10.                                             \tag{I22.3}
\]
No dominance or new determinant class for (I22.3) is asserted.

\paragraph{1. Primitive projective groups and the actual field obstruction.}
Write $Z=\Gamma\cap\mathbf C^*I$ and
$\bar\Gamma=\Gamma/Z\subset\operatorname{PGL}_3(\mathbf C)$.
Every scalar in $Z$ belongs to $K$ and has finite order.
The roots of unity in our fields are $\mu_4$ for $K=\mathbf Q(i)$
and $\{\pm1\}$ otherwise. In particular their orders are
powers of two and they contain no nontrivial cube root of unity.

The classification input is Dolgachev--Iskovskikh,
\emph{Finite subgroups of the plane Cremona group},
Theorem 4.8: the primitive finite projective groups are
$A_5$, $A_6$, $\operatorname{PSL}_2(7)$, and the three
Hessian groups of orders $36,72,216$. Each Hessian group
contains $C_3^2$ as a normal subgroup. The published chapter
is in \emph{Algebra, Arithmetic, and Geometry}, Progress
in Mathematics 269 (2009), 443--548,
DOI \texttt{10.1007/978-0-8176-4745-2\_11}.
The author-hosted typeset proof places Theorem 4.8 on its
printed p.459, PDF page 23. That is the exact version and
locator checked here; its preceding Theorem 4.7 and
imprimitive argument are on pp.458--459.

The cases $A_5,A_6$ supply a projective element of order five.
For an actual lift $g\in\Gamma$ of finite order $n$, the
matrix $g^{n/5}$ has exact order five. A nonidentity
eigenvalue is primitive of order five. Since the unique
quadratic subfield of $\mathbf Q(\zeta_5)$ is real,
$K\cap\mathbf Q(\zeta_5)=\mathbf Q$, so its minimal
polynomial over $K$ has degree four. It cannot divide the
degree-three characteristic polynomial of this actual matrix.
This excludes both groups, without substituting a character
table for the representation.
Explicitly, the degree-four cyclotomic Galois group is cyclic,
and its unique quadratic fixed field is generated by the real
number $\zeta_5+\zeta_5^{-1}$, which satisfies $x^2+x-1=0$.

For a Hessian group, take the full inverse image $B$ of
its displayed $C_3^2$. The scalar kernel has order prime
to three, so a Sylow-three subgroup $D\subset B$ maps
isomorphically onto $C_3^2$: its order is nine and its
intersection with the scalar kernel is trivial. This is
an actual faithful subgroup of $\operatorname{GL}_3(K)$.
Over $K(\omega)$, its commuting finite-order matrices
decompose into characters. The nontrivial characters occur
in conjugate pairs $\chi,\chi^{-1}$, since
$[K(\omega):K]=2$ and the matrices are over $K$.
In dimension three there is at most one such pair.
The two characters have the same order-three kernel;
the remaining trivial character does not remove it.
The resulting representation of $C_3^2$ is not faithful,
a contradiction. This excludes all three Hessian groups.

For the remaining primitive projective group
$\operatorname{PSL}_2(7)$, an order-seven lift similarly
gives an actual matrix of order seven. Its primitive
eigenvalue has degree at most three over $K$.
The cyclotomic field $\mathbf Q(\zeta_7)$ has degree six
and unique quadratic subfield $K_7$, so necessarily $K=K_7$.
For completeness, with
$\lambda=\zeta_7+\zeta_7^2+\zeta_7^4$ one has
$\lambda+\bar\lambda=-1$, $\lambda\bar\lambda=2$,
and hence $\lambda^2+\lambda+2=0$. This identifies that
quadratic subfield by an explicit element.

The projective group $\operatorname{PSL}_2(7)$ is perfect.
An elementary matrix proof suffices: the upper and lower
unipotent matrices generate $\operatorname{SL}_2(\mathbf F_7)$;
conjugation by $\operatorname{diag}(2,2^{-1})$ multiplies
the upper unipotent parameter by $4$, so every such
unipotent is a commutator, as $4-1$ is invertible.
The same holds for lower unipotents, and the quotient
by $\{\pm I\}$ is perfect as well.
Now use the actual determinant, not a choice of projective lifts.
Set $\Gamma_0=\ker(\det:\Gamma\to\{\pm1\})$.
If $Z=1$, the determinant factors through the perfect
projective group and is trivial. If $Z=\{\pm I\}$,
each projective element has a determinant-one lift by
multiplication by $-I$, whose determinant in degree three
is $-1$. In both cases the projective map restricts to
an isomorphism
\[
 \Gamma_0\xrightarrow{\ \sim\ }\operatorname{PSL}_2(7),
 \qquad \Gamma=\Gamma_0\ \hbox{or}\ \Gamma_0\times\{\pm I\}.
                                                               \tag{I22.4}
\]

Here is an explicit $K_7$ model for $\Gamma_0$.
Use the three involutions of HD-S13:
\[
 s_1=\begin{pmatrix}-1&1&\lambda\\0&1&0\\0&0&1\end{pmatrix},\quad
 s_2=\begin{pmatrix}1&0&0\\1&-1&1\\0&0&1\end{pmatrix},\quad
 s_3=\begin{pmatrix}1&0&0\\0&1&0\\-1-\lambda&1&-1\end{pmatrix},
 \qquad L=\langle s_1s_2,s_2s_3\rangle.                       \tag{I22.5}
\]
The exact matrix certificate enumerates $168$ elements
of $L$, verifies that its commutator subgroup is all of $L$,
and computes character norm $168$. Thus it is absolutely
irreducible and perfect. Its scalar center is trivial,
as follows also from determinant one and $\mu_3(K_7)=1$.
A perfect irreducible degree-three group cannot be
imprimitive: its permutation image on three lines would
be a perfect subgroup of $S_3$, hence trivial, which
would make its matrices diagonal. Its projective group
is therefore the order-$168$ case of Theorem 4.8.
It preserves the positive Hermitian matrix
\[
 H_L=\begin{pmatrix}2&-1&-\lambda\\-1&2&-1\\1+\lambda&-1&2\end{pmatrix},
 \qquad \text{leading minors }2,3,1.                         \tag{I22.6}
\]
The prior root-data morphism to Cohen's published
$(4.15)$, p.407, is the actual matrix $T$ in HD-S13;
no new root-data identification is presumed here.

The conjugacy in Theorem 4.8 is projective. To obtain the
claimed actual $K$-linear map, take its abstract group
identification $\phi:\Gamma_0\to L$ and a complex matrix
$P$ giving the projective conjugacy. The scalar factors in
$P\rho(g)P^{-1}=c_g\rho_L(\phi(g))$ form a character
of the perfect group $\Gamma_0$, so every $c_g=1$.
The finite intertwiner equations
$U\rho(g)=\rho_L(\phi(g))U$ have coefficients in $K$
and a nonzero complex solution. Their nullspace therefore
has a nonzero $K$-solution; any nonzero intertwiner is
invertible by absolute irreducibility. This is the exact
$K$-linear conjugation used in (I22.2).

\paragraph{2. An imprimitive line system and its precise Galois descent.}
Suppose the absolutely irreducible group is complex-imprimitive.
Choose its three permuted complex lines, let
$P\subset S_3$ be the transitive permutation image, and
let $D$ be the actual kernel in $\Gamma$. Its matrices
are diagonal on these lines and $P=C_3$ or $S_3$.

First $D$ is not scalar. Otherwise take a lift $T$ of a
three-cycle. Its cube is a scalar $c\in\mu(K)$.
The cube map on $\mu(K)$ is bijective, so rescale $T$
by a scalar of $K$ to have $T^3=I$. Its three-cycle
monomial matrix has characteristic polynomial $x^3-1$;
therefore $\ker(T-I)$ is a one-dimensional $K$-space.
Every element of $\Gamma$ conjugates $T$ to a scalar
times $T$ or $T^{-1}$, because the three-cycle subgroup
is normal in $P$. That scalar has cube one and belongs
to $K$, so it equals one. Consequently this $K$-line
is $\Gamma$-invariant, contradicting absolute irreducibility.

The three characters $\chi_1,\chi_2,\chi_3$ of $D$
are now distinct. Indeed equality defines a $P$-invariant
equivalence relation on three points; transitivity makes
it either discrete or universal, and the universal case
would make $D$ scalar. Thus the chosen lines are exactly
the three common eigenspaces of the actual subgroup $D$.
Galois fixes every matrix of $D$ and of $\Gamma$, so it
permutes these eigenspaces and its permutation commutes
with $P$. If $P=S_3$, its centralizer in $S_3$ is trivial,
and all three lines descend to $K$. If $P=C_3$, that
centralizer is $C_3$, so either all three lines descend
or they form one Galois orbit with cyclic cubic field
of definition $E/K$.

In the latter case, choose one character $\chi$.
The field of its values is exactly $E$, since its Galois
stabilizer is the stabilizer of the corresponding unique
eigenspace. The other two characters are its conjugates.
Their common kernel is trivial, so already $\chi$ is
faithful: an element killed by $\chi$ is killed by all
its conjugates and hence acts identically on all three
lines. Therefore $D$ is cyclic, its image is $\mu_m$,
and
\[
                 E=K(\zeta_m),\qquad [E:K]=3.              \tag{I22.7}
\]
The degree formula gives $\varphi(m)=3$ or $6$; the
first is impossible. Solving $\varphi(m)=6$ gives
$m=7,9,14,18$: every prime divisor has $p-1\mid6$,
so only $2,3,7$ occur, and the prime-power factors give
exactly these four integers. For $m=9,18$ the unique
quadratic subfield is $\mathbf Q(\omega)$, outside our
list. Thus
\[
                  K=K_7,\qquad E=\mathbf Q(\zeta_7),
                         \qquad m=7\ \hbox{or}\ 14.       \tag{I22.8}
\]

The actual projectors, rather than just the permutation
of formal characters, are
\[
 p_j=\frac1{|D|}\sum_{d\in D}\chi_j(d)^{-1}\rho(d)
                   \in M_3(E).                            \tag{I22.9}
\]
They are the three rank-one pairwise orthogonal idempotents
over $E$ and sum to $I$. Their Galois permutations prove
the field-of-definition assertion just used. Equivalently,
a generator of $D$ has an irreducible degree-three minimal
polynomial over $K$, and the algebra it generates is an
actual embedded field $E\subset\operatorname{End}_K(M)$.
Thus $M$ has rank one over $E$.

This also determines the exact normalizer of $D$ over $K$.
An element normalizing $D$ normalizes its generated field,
and hence induces a $K$-automorphism of $E$. Its centralizer
is $E^*$ because $M$ has $E$-rank one. In an $E$-basis,
every normalizer element therefore acts by
\[
                   x\longmapsto a\sigma^j(x),\qquad
      a\in E^*,\quad j\in\{0,1,2\},\quad
      \sigma(\zeta_7)=\zeta_7^2.                           \tag{I22.10}
\]
Conversely these maps normalize the multiplication group
$\mu_m$. Hence this is equality with the actual normalizer,
not just an embedding in a conjectured semilinear model.

Since $\Gamma/D=C_3$, choose $T=a\sigma$ after reversing
the generator when necessary. Then
\[
              T^3=N_{E/K}(a)=c\in D\cap K^*\subseteq\{\pm1\}.
                                                               \tag{I22.11}
\]
Multiplication of $T$ by $c^{-1}\in D$ gives
$T'=a'\sigma$ with $N_{E/K}(a')=1$ and $(T')^3=I$.
There is a fully explicit basis correction. The monomial
three-cycle $T'$ has characteristic polynomial $x^3-1$,
so $1+T'+(T')^2$ is nonzero. Choose
$z\in\{1,\zeta_7,\zeta_7^2\}$ for which
\[
 b=z+a'\sigma(z)+a'\sigma(a')\sigma^2(z)\ne0.
                                                               \tag{I22.12}
\]
The norm-one equation gives $a'\sigma(b)=b$.
The actual $K$-linear basis change $x\mapsto bx$ then
conjugates $T'$ to $\sigma$ and leaves all multiplication
operators of $D$ unchanged. Thus $\Gamma$ is literally
the semidirect product in (I22.2), with or without the
central $-I$. This proves the splitting of the finite
extension from its matrices and norm equation.

For reference, in the basis $1,\zeta_7,\zeta_7^2$, choosing
the generator with trace $\lambda$, the two matrices are
\[
 A=\begin{pmatrix}0&0&1\\1&0&1+\lambda\\0&1&\lambda\end{pmatrix},
 \qquad
 S=\begin{pmatrix}1&0&\lambda\\0&0&-1\\0&1&-1\end{pmatrix},
 \quad A^7=S^3=I,\quad SAS^{-1}=A^2.                       \tag{I22.13}
\]
Their group has $21$ elements; adjoining $-I$ gives $42$.
They preserve the actual trace Hermitian form
\[
 H_{21}=\begin{pmatrix}3&\lambda&\lambda\\
 \bar\lambda&3&\lambda\\\bar\lambda&\bar\lambda&3\end{pmatrix},
 \qquad \text{leading minors }3,7,7,                        \tag{I22.14}
\]
which is $\operatorname{Tr}_{E/K}(\bar x y)$ on $E$.
The supplementary author certificate verifies these displayed
matrix identities by exact arithmetic.

If the line system instead descends to $K$, its stabilizer
characters have values in $\mu(K)$. The character of the
induced degree-three representation is a sum of such values.
For $K\ne\mathbf Q(i)$ in our list it is therefore rational,
contradicting the stipulated exact nonreal character field.
For $K=\mathbf Q(i)$ this is precisely the allowed first
case of the statement. Together with paragraph1 this proves
the complete field and group list.

\paragraph{3. The actual S45 source map and the spectral exclusion.}
Retain the S45 cover with its actual product group
$H=S_s\times\Gamma$, $r=s-1\ge9$, and the coefficient
$W\otimes_K M$, where $W=\operatorname{Aug}(K^s)$.
Changing the coefficient basis by any of the intertwiners
above relabels the same deck group. First use the
integer-cleared matrix-coefficient corner morphism of
HD-S45, paragraph1, to carry the originally prescribed
line in $W\otimes_KM$ to $Kw_0\otimes_Ku$, where
$w_0=(1,\ldots,1,1-s)$ spans the fixed line of $S_{s-1}$
in $W$. That established polarized isogeny is the one
retaining the original factor and its determinant class.
An intertwiner $U$ on $M$ then acts on the full coefficient
space by the actual matrix $I_W\otimes U$.

To check the orthogonal projector explicitly, write the
actual coefficient form as $h_W\otimes h$ and the model
form as $h_W\otimes H$. Schur's lemma and conjugation give
$h=cU^*HU$ with $c\in\mathbf Q_{>0}$. The factor $c$
cancels in $u(u^*hu)^{-1}u^*h$. Thus the Gamma corner
$e_u$ on $M$ has model matrix $Ue_uU^{-1}$, and the full
corner is transported by $I_W\otimes U$ to the orthogonal
projector onto $w_0\otimes Uu$. These are the displayed
coefficient maps, after the preceding actual corner isogeny.
Explicitly, for the actual connected cover $C$ put
\[
 \begin{gathered}
 L_0=S_{s-1},\quad X=C/(L_0\times1),\quad
 Z_0=C/(S_s\times1),\quad Y=C/(L_0\times\Gamma),\\
 \alpha:C\to X,\qquad f:X\to Z_0,\qquad\deg f=s.
 \end{gathered}                                             \tag{I22.15}
\]
For the orthogonal primitive rational idempotent $e_u$
of the conjugate character pair, the transported sixfold factor is
the connected image of an integer multiple of
\[
 \pi=e_u\left(1-s^{-1}f^*f_*\right),\qquad
 |L_0|^{-1}\alpha^*\pi\alpha_*
       =(\operatorname{av}_{L_0}-\operatorname{av}_{S_s})e_u.
                                                               \tag{I22.16}
\]
The norm identities give the stated equality and the
polarization multiplier $|L_0|$ for $\alpha^*$.
The diagonal map $v\mapsto(v,\ldots,v)$ and the average
$s^{-1}\sum v_i$ split the parabolic cohomology on $Y$
into the base $\Gamma$-factor and this orthogonal complement;
on a ramified cycle of length $\ell$ they use exactly
$\gamma^\ell$. These are the same actual maps (S45.10)--(S45.12).
Every coefficient conjugation above transports the invariant
form and the specified line, so it leaves this corner identity
and its polarized factor unchanged.

For the two groups generated by $A,S$ in (I22.13), every
nonscalar element has three distinct complex eigenvalues.
For $A^a$, $a\not\equiv0\pmod7$, the eigenvalues are
$\zeta_7^a,\zeta_7^{2a},\zeta_7^{4a}$. For $A^aS$
or $A^aS^2$, the cube is $I$ and the permutation of the
three eigenspaces is a three-cycle, so the characteristic
polynomial is $x^3-1$. Multiplication by the scalar $-I$
preserves distinctness. If $B$ is any such nonscalar
matrix and $P_\rho$ any permutation matrix on $W$, each
eigenvalue of $W$ matches at most one eigenline of $B$.
Consequently
\[
 \dim\ker(P_\rho\otimes B-I)\le r,
 \qquad \operatorname{codim}(P_\rho\otimes B-I)\ge2r.
                                                               \tag{I22.17}
\]

The actual projective $\Gamma$-tuple has at least three
nonidentity entries. Otherwise the product-one relation
makes it cyclic; a finite group with cyclic projective
image has matrices that are scalar multiples of powers
of one semisimple matrix and is reducible in degree three.
All remaining nonidentity inertia in the current S45
reduction has codimension at least three. Its established
Euler and dimension conditions are
$b\ge12$ and $\sum c_j=6r+6$. Formula (I22.17) therefore
gives the contradiction
\[
             6r+6=\sum c_j\ge6r+3(b-3)\ge6r+27.           \tag{I22.18}
\]
This excludes the two nonsplit-line monomial groups as
actual rank-nine tensor sources, leaving only (I22.3).

For a $K$-defined line system the induced representation,
quotient and polarization maps are (S45.13), with faithful
block image $S_s\times\chi(J)$. Exact field forces
$\chi(J)=\mu_4$ in the Gaussian case, so the already proved
scalar-augmentation source conclusion applies through that
actual block. Its determinant classes are $[-1]$ and $[-3]$.
For the other six fields outside $K_7$ the exact-field
representation is absent by the classification above.
This is the stated closure of the S45 multiplicity-three
branch over those seven fields. It makes no assertion about
additional endomorphisms on special subfamilies of rational
coefficient factors, or about coefficient models outside
the precise S45 reduction.
